%% chapters/assignment/ch5_so.tex
%% 第5章　システム最適配分（SO）と混雑課金

\section{システム最適配分（SO）と混雑課金}
\label{sec:ch5-so}

\sref{sec:ch3-ue}で扱った利用者均衡配分（UE）では，各利用者が自己の旅行時間を最小化するように
経路選択を行う結果として均衡が成立することを示した（\ssref{subsec:beckmann}）．
しかし，\sref{sec:ch3-ue}の Braess のパラドックス（\ssref{subsec:braess}）が示すように，
個人の合理的行動の集積が必ずしも社会全体にとって望ましい結果をもたらすとは限らない．
個々の利用者は道路に参入するとき，自分の旅行時間しか考慮しないが，
実際には他の利用者の旅行時間も増大させている——これが\textbf{混雑の外部性}である．

本節では，この外部性を明示的に扱い，ネットワーク全体の\textbf{総旅行時間を最小化する}
配分として定義される\textbf{システム最適配分}（System Optimal, SO）を導入する（\ssref{subsec:wardrop2}）．
続いて UE と SO の相違を比較し（\ssref{subsec:ue-so-comparison}），
混雑課金（Pigouvian Tax）によって UE を SO と一致させる
\textbf{限界費用課金}の考え方を説明する（\ssref{subsec:congestion-toll}）．
最後に，余剰分析の手法を用いて，
UE における死荷重損失と混雑課金の効果を定量的に整理する（\ssref{subsec:surplus-analysis}）\cite{Pigou1920,Beckmann1956}．

% ---------------------------------------------------------------
\subsection{Wardropの第2原則と定式化}
\label{subsec:wardrop2}
% ---------------------------------------------------------------

利用者均衡配分の根拠となる Wardrop の第1原則に対応して，
\textbf{Wardrop の第2原則}（Wardrop's Second Principle）は次のように述べる\cite{Wardrop1952}．

\begin{quotation}
  \noindent
  「ネットワーク全体の総旅行時間（旅行費用）が最小となるように交通量が配分される．」
\end{quotation}

この原則を実現する配分が\textbf{システム最適配分}（SO）である．
SO は，中央調整者（交通管理機関など）がすべての利用者の経路を
社会的観点から最適に割り当てることができる場合の理想的な配分に対応する．

\subsubsection*{SO-Primal 問題の定式化}

SO は，制約条件を UE/FD-Primal（\ssref{subsec:beckmann}参照）と共有しつつ，
\textbf{目的関数だけが異なる}最適化問題として定式化される\cite{Beckmann1956}．

\begin{description}
  \item[SO-Primal]
    \begin{equation}
      \min_{\boldsymbol{x},\,\boldsymbol{f}}\;
      Z_{SO}(\boldsymbol{x})
      = \sum_{a \in A} x_a\, t_a(x_a)
      \label{eq:so-primal}
    \end{equation}
    \begin{align*}
      \text{s.t.} \quad
      & x_a = \sum_{rs \in \Omega}\sum_{k \in K_{rs}} f_k^{rs}\,\delta_{a,k}^{rs}
        \qquad \forall a \in A \\
      & \sum_{k \in K_{rs}} f_k^{rs} = q_{rs}
        \qquad \forall rs \in \Omega \\
      & f_k^{rs} \geq 0
        \qquad \forall k \in K_{rs},\; \forall rs \in \Omega
    \end{align*}
\end{description}

目的関数 $Z_{SO}(\boldsymbol{x}) = \sum_{a} x_a t_a(x_a)$ は
ネットワーク全体の\textbf{総旅行時間}（各リンクの旅行時間 $t_a$ と交通量 $x_a$ の積の和）を表す．
リンクパフォーマンス関数 $t_a(x_a)$ が単調非減少かつ凸であれば
$Z_{SO}$ も凸関数となり，大域最適解が一意に定まる．

\subsubsection*{Frank-Wolfe 法の適用}

SO を解く方法は UE の場合（\ssref{subsec:frank-wolfe}参照）と同様に
Frank-Wolfe 法をそのまま適用できる\cite{FrankWolfe1956}．
ただし，線形化ステップ（All-or-Nothing 配分）で用いるリンクコストが異なる点に注意する．

UE-Frank-Wolfe では各リンクの旅行時間 $t_a(x_a^{(n)})$（私的限界費用）を
最短経路探索のコストとして用いる\cite{土木学会1998}．
一方，SO-Frank-Wolfe では目的関数の勾配，すなわち
\begin{equation}
  \frac{\partial Z_{SO}}{\partial x_a}
  = \frac{d}{dx_a}\bigl[x_a\, t_a(x_a)\bigr]
  = t_a(x_a) + x_a\, t_a'(x_a)
  \label{eq:smc-gradient}
\end{equation}
を各リンクのコストとして用いる．
式\eqref{eq:smc-gradient}の右辺は後述の\textbf{社会的限界費用}（SMC）に一致しており，
SO の KKT 条件は「全使用経路の社会的限界費用の和が等しい」という
Wardrop の第2原則に対応する．

% ---------------------------------------------------------------
\subsection{UEとSOの比較}
\label{subsec:ue-so-comparison}
% ---------------------------------------------------------------

UE と SO の本質的な違いを一言で表せば，
「個人の最適化」対「全体の最適化」である．

UE では各利用者が\textbf{私的限界費用}（PMC）——自分が経験する旅行時間 $t_a(x_a)$——
のみを考慮して経路を選択するため，
自分が道路に加わることで他の利用者に与える追加的な旅行時間遅延（外部コスト）は無視される．
この外部性の無視が，UE 配分を社会的に非効率にする根本原因である．

SO では，すべての利用者の旅行時間を考慮することで，
外部コストを含む\textbf{社会的限界費用}（SMC）が等化される均衡が実現し，
ネットワークが全体として最も効率的に利用される．

\begin{table}[hbtp]
  \centering
  \caption{利用者均衡配分（UE）とシステム最適配分（SO）の比較}
  \label{tab:ue-so-comparison}
  \begin{tabular}{lll}
    \toprule
    比較項目 & UE & SO \\
    \midrule
    原則 & Wardrop の第1原則 & Wardrop の第2原則 \\
    目的 & 個人の旅行時間最小化 & 社会全体の総旅行時間最小化 \\
    目的関数 &
      $\displaystyle\sum_{a}\int_0^{x_a} t_a(w)\,dw$ &
      $\displaystyle\sum_{a} x_a\, t_a(x_a)$ \\[4pt]
    最短経路コスト & $t_a(x_a)$（PMC） & $t_a(x_a)+x_a t_a'(x_a)$（SMC） \\
    均衡交通量 & $x_{UE}$（大） & $x_{SO}$（小） \\
    均衡旅行時間 & $t_{UE} = t_a(x_{UE})$（長） & $t_{SO} = t_a(x_{SO})$（短） \\
    外部コスト & 無視 & 内部化 \\
    \bottomrule
  \end{tabular}
\end{table}

一般に $x_{SO} \leq x_{UE}$ および $t_{SO} \leq t_{UE}$ が成り立つ
（需要が価格に反応する場合は厳密不等号）．
UE では各利用者が外部コストを負担しないまま道路を利用するため，
社会的に最適な水準 $x_{SO}$ を超えた「過剰な交通量」が発生する．

\subsubsection*{Braess のパラドックスとの関係}

\sref{sec:ch3-ue}（\ssref{subsec:braess}）で示した Braess のパラドックスは，
UE の非効率性の顕著な例でもある．
道路を追加したとき，UE 配分では総旅行時間がかえって悪化するにもかかわらず，
SO 配分においてこのような逆説は生じない：最適化が社会全体の視点から行われるため，
新設リンクが実際に有益な場合にのみそれが積極的に利用される\cite{Braess1968}．

% ---------------------------------------------------------------
\subsection{混雑課金（限界費用課金）}
\label{subsec:congestion-toll}
% ---------------------------------------------------------------

UE を SO に一致させるための政策的手段が\textbf{混雑課金}（Congestion Toll），
あるいは\textbf{限界費用課金}（Marginal Cost Pricing）である．
この考えは Pigou（1920）に遡り，外部性を価格メカニズムを通じて内部化する
\textbf{ピグー税}（Pigouvian Tax）の交通版に相当する\cite{Pigou1920}．

\subsubsection*{社会的限界費用の分解}

1 台の車両が混雑したリンク $a$ に追加されるとき，
ネットワーク全体（リンク $a$ のみを考える単純な場合）に与えるコスト変化は次のように分解される：
\begin{align}
  \text{SMC}_a(x_a)
  &= \frac{d}{dx_a}\bigl[x_a\, t_a(x_a)\bigr]
   = t_a(x_a)
     + x_a\, \frac{d\, t_a(x_a)}{dx_a}  \notag \\
  &= \underbrace{t_a(x_a)}_{\text{私的限界費用（PMC）}}
     + \underbrace{x_a\, t_a'(x_a)}_{\text{外部コスト}}
  \label{eq:smcdecomp}
\end{align}
第1項の私的限界費用（PMC）は，利用者自身が経験する旅行時間である．
第2項の外部コストは，1 台が追加されることで既存の $x_a$ 台すべての旅行時間が
$t_a'(x_a)$ だけ増加することによる社会的損失を表す．

UE では利用者が外部コスト $x_a t_a'(x_a)$ を無視して意思決定するため，
個人均衡（UE）において過剰な交通量が発生する．

\subsubsection*{最適課金額の導出}

各リンク $a$ に課金額 $\tau_a$ を設定すると，利用者が直面するコストは
$t_a(x_a) + \tau_a$ となる．
このとき，利用者の私的最適化が社会的最適解 $\boldsymbol{x}^{SO}$ を実現するためには，
SO 解における社会的限界費用を課金後の私的限界費用と一致させればよい\cite{Beckmann1956}：
\begin{equation}
  \tau_a^* = x_a^{SO}\, t_a'\bigl(x_a^{SO}\bigr)
  \label{eq:optimal-toll}
\end{equation}
この最適料金 $\tau_a^*$ を各リンクに設定したとき，
修正後の私的費用に基づく UE 解は SO 解と一致する．
すなわち，$x_{UE}^{\text{toll}} = x_{SO}$ が成立する．

\subsubsection*{BPR 関数への適用}

BPR 関数 $t_a(x_a) = t_{a0}\left\{1 + \alpha(x_a/C_a)^\beta\right\}$（\ssref{subsec:bpr}）
に対しては：
\begin{equation}
  t_a'(x_a) = t_{a0}\,\alpha\,\beta\,\frac{x_a^{\beta-1}}{C_a^{\beta}}
  \label{eq:bpr-deriv}
\end{equation}
したがって最適課金額は：
\begin{equation}
  \tau_a^* = x_a^{SO} \cdot t_{a0}\,\alpha\,\beta\,
               \frac{(x_a^{SO})^{\beta-1}}{C_a^{\beta}}
           = t_{a0}\,\alpha\,\beta\,\frac{(x_a^{SO})^{\beta}}{C_a^{\beta}}
  \label{eq:bpr-toll}
\end{equation}
となる．混雑が激しく $x_a^{SO}/C_a$ が大きいリンクほど高い料金が設定されることがわかる．

% ---------------------------------------------------------------
\subsection{余剰分析}
\label{subsec:surplus-analysis}
% ---------------------------------------------------------------

混雑課金の効果を定量的に評価するために，
ミクロ経済学の\textbf{余剰分析}（Welfare Analysis）の枠組みを用いる\cite{土木学会1998}．
ここでは，交通量の変化に伴う社会的厚生の変化を，
需要曲線・PMC 曲線・SMC 曲線の関係として図示する（図\ref{fig:surplus-analysis}）．

\begin{figure}[hbtp]
  \centering
  %% fig_surplus_analysis.tex
%% 余剰分析：UE と SO の比較，最適混雑課金による死荷重損失の解消
%% Usage: %% fig_surplus_analysis.tex
%% 余剰分析：UE と SO の比較，最適混雑課金による死荷重損失の解消
%% Usage: %% fig_surplus_analysis.tex
%% 余剰分析：UE と SO の比較，最適混雑課金による死荷重損失の解消
%% Usage: \input{assets/assignment/fig_surplus_analysis.tex}
%%
%% 模式図のパラメータ（すべて正規化された単位）
%%   需要曲線     D(x)   = 4 - 0.5x         （右下がり）
%%   私的限界費用  PMC(x) = 1 + 0.5x = t_a(x_a)
%%   社会的限界費用 SMC(x) = 1 + x   = d/dx[x * t_a(x_a)]
%%
%%   SO 点 (D = SMC): x_SO = 2,  t_SO = PMC(2) = 2,  価格(含む課金) = 3
%%   UE 点 (D = PMC): x_UE = 3,  t_UE = PMC(3) = 2.5
%%   最適課金 τ* = SMC(x_SO) - PMC(x_SO) = 3 - 2 = 1
%%   自由流旅行時間 t_0 = PMC(0) = SMC(0) = 1
%%   死荷重損失の三角形 頂点: SO点(2,3), UE点(3,2.5), SMC(x_UE)=(3,4)

\begin{tikzpicture}[
  every node/.style={font=\small},
  xscale=1.6,    %% 横方向の拡大率（x 軸 0-4）
  yscale=1.4     %% 縦方向の拡大率（y 軸 0-4.5）
]

%% =============================================================
%% 塗り潰し（背面レイヤー）
%% =============================================================

%% 死荷重損失「紫の三角形」:  SO点 → SMC(x_UE) → UE点 で囲まれる三角形
%% 境界: 上辺 = SMC 曲線 (2,3)→(3,4), 右辺 = 垂直線 (3,4)→(3,2.5),
%%       下辺 = D 曲線 (3,2.5)→(2,3)
\fill[violet!30] (2, 3) -- (3, 4) -- (3, 2.5) -- cycle;

%% =============================================================
%% 曲線（3 本）
%% =============================================================

%% 需要曲線  D(x) = 4 - 0.5x
\draw[very thick, blue!75!black]
  (0, 4) -- (3.6, 2.2)
  node[right] {$D$};

%% 私的限界費用  PMC(x) = t_a(x_a) = 1 + 0.5x
\draw[very thick, green!55!black]
  (0, 1) -- (3.6, 2.8)
  node[right] {$\text{PMC} = t_a(x_a)$};

%% 社会的限界費用  SMC(x) = 1 + x （= d/dx[x * t_a] = t_a + x_a t_a'）
\draw[very thick, red!65!black]
  (0, 1) -- (3.3, 4.3)
  node[right, align=left] {$\text{SMC} = t_a + x_a\, t_a'$};

%% =============================================================
%% 補助破線
%% =============================================================

%% x_SO の垂直破線
\draw[gray, dashed, thin] (2, 0) -- (2, 3);

%% x_UE の垂直破線
\draw[gray, dashed, thin] (3, 0) -- (3, 2.5);

%% t_SO = PMC(x_SO) = 2 の水平破線
\draw[gray, dashed, thin] (0, 2) -- (2, 2);

%% t_UE = PMC(x_UE) = 2.5 の水平破線
\draw[gray, dashed, thin] (0, 2.5) -- (3, 2.5);

%% =============================================================
%% 均衡点（塗り潰し丸）
%% =============================================================

%% SO 点（D と SMC の交点, x=2, y=3）
\fill (2, 3) circle[radius=2.3pt];

%% UE 点（D と PMC の交点, x=3, y=2.5）
\fill (3, 2.5) circle[radius=2.3pt];

%% PMC(x_SO) の点（課金の下端, x=2, y=2）
\fill[gray!60] (2, 2) circle[radius=1.8pt];

%% =============================================================
%% 軸（前面レイヤー）
%% =============================================================

\draw[thick, ->] (-0.15, 0) -- (4.0, 0) node[right] {$x_a$（交通量）};
\draw[thick, ->] (0, -0.15) -- (0, 4.7) node[above] {費用（旅行時間）};

%% 原点ラベル
\node[below left, font=\small] at (0, 0) {O};

%% =============================================================
%% 軸の目盛り・ラベル
%% =============================================================

%% x 軸
\draw[thin] (2, 0) -- (2, -0.1);
\draw[thin] (3, 0) -- (3, -0.1);
\node[below]          at (2, 0)   {$x_{SO}$};
\node[below]          at (3, 0)   {$x_{UE}$};

%% y 軸
\draw[thin] (-0.1, 1)   -- (0, 1);
\draw[thin] (-0.1, 2)   -- (0, 2);
\draw[thin] (-0.1, 2.5) -- (0, 2.5);

\node[left, font=\footnotesize] at (0, 1)   {$t_0$};
\node[left]                     at (0, 2)   {$t_{SO}$};
\node[left]                     at (0, 2.5) {$t_{UE}$};

%% =============================================================
%% τ* 注釈（最適課金額：PMC(x_SO)〜SMC(x_SO) の垂直両矢印）
%% =============================================================

\draw[<->, >=stealth, thick]
  (2.2, 2.0) -- (2.2, 3.0)
  node[midway, right, font=\footnotesize] {$\tau^*$};

%% =============================================================
%% 領域ラベル
%% =============================================================

%% 死荷重損失ラベル（三角形の内側）
\node[font=\footnotesize, align=center, violet!70!black]
  at (2.72, 3.15) {死荷重\\[-2pt]損失};

%% =============================================================
%% 均衡点のラベル
%% =============================================================

\node[above left, font=\footnotesize] at (2.2, 3.1)    {SO点};
\node[below right, font=\footnotesize] at (2.8, 2.4) {UE点};

\end{tikzpicture}

%%
%% 模式図のパラメータ（すべて正規化された単位）
%%   需要曲線     D(x)   = 4 - 0.5x         （右下がり）
%%   私的限界費用  PMC(x) = 1 + 0.5x = t_a(x_a)
%%   社会的限界費用 SMC(x) = 1 + x   = d/dx[x * t_a(x_a)]
%%
%%   SO 点 (D = SMC): x_SO = 2,  t_SO = PMC(2) = 2,  価格(含む課金) = 3
%%   UE 点 (D = PMC): x_UE = 3,  t_UE = PMC(3) = 2.5
%%   最適課金 τ* = SMC(x_SO) - PMC(x_SO) = 3 - 2 = 1
%%   自由流旅行時間 t_0 = PMC(0) = SMC(0) = 1
%%   死荷重損失の三角形 頂点: SO点(2,3), UE点(3,2.5), SMC(x_UE)=(3,4)

\begin{tikzpicture}[
  every node/.style={font=\small},
  xscale=1.6,    %% 横方向の拡大率（x 軸 0-4）
  yscale=1.4     %% 縦方向の拡大率（y 軸 0-4.5）
]

%% =============================================================
%% 塗り潰し（背面レイヤー）
%% =============================================================

%% 死荷重損失「紫の三角形」:  SO点 → SMC(x_UE) → UE点 で囲まれる三角形
%% 境界: 上辺 = SMC 曲線 (2,3)→(3,4), 右辺 = 垂直線 (3,4)→(3,2.5),
%%       下辺 = D 曲線 (3,2.5)→(2,3)
\fill[violet!30] (2, 3) -- (3, 4) -- (3, 2.5) -- cycle;

%% =============================================================
%% 曲線（3 本）
%% =============================================================

%% 需要曲線  D(x) = 4 - 0.5x
\draw[very thick, blue!75!black]
  (0, 4) -- (3.6, 2.2)
  node[right] {$D$};

%% 私的限界費用  PMC(x) = t_a(x_a) = 1 + 0.5x
\draw[very thick, green!55!black]
  (0, 1) -- (3.6, 2.8)
  node[right] {$\text{PMC} = t_a(x_a)$};

%% 社会的限界費用  SMC(x) = 1 + x （= d/dx[x * t_a] = t_a + x_a t_a'）
\draw[very thick, red!65!black]
  (0, 1) -- (3.3, 4.3)
  node[right, align=left] {$\text{SMC} = t_a + x_a\, t_a'$};

%% =============================================================
%% 補助破線
%% =============================================================

%% x_SO の垂直破線
\draw[gray, dashed, thin] (2, 0) -- (2, 3);

%% x_UE の垂直破線
\draw[gray, dashed, thin] (3, 0) -- (3, 2.5);

%% t_SO = PMC(x_SO) = 2 の水平破線
\draw[gray, dashed, thin] (0, 2) -- (2, 2);

%% t_UE = PMC(x_UE) = 2.5 の水平破線
\draw[gray, dashed, thin] (0, 2.5) -- (3, 2.5);

%% =============================================================
%% 均衡点（塗り潰し丸）
%% =============================================================

%% SO 点（D と SMC の交点, x=2, y=3）
\fill (2, 3) circle[radius=2.3pt];

%% UE 点（D と PMC の交点, x=3, y=2.5）
\fill (3, 2.5) circle[radius=2.3pt];

%% PMC(x_SO) の点（課金の下端, x=2, y=2）
\fill[gray!60] (2, 2) circle[radius=1.8pt];

%% =============================================================
%% 軸（前面レイヤー）
%% =============================================================

\draw[thick, ->] (-0.15, 0) -- (4.0, 0) node[right] {$x_a$（交通量）};
\draw[thick, ->] (0, -0.15) -- (0, 4.7) node[above] {費用（旅行時間）};

%% 原点ラベル
\node[below left, font=\small] at (0, 0) {O};

%% =============================================================
%% 軸の目盛り・ラベル
%% =============================================================

%% x 軸
\draw[thin] (2, 0) -- (2, -0.1);
\draw[thin] (3, 0) -- (3, -0.1);
\node[below]          at (2, 0)   {$x_{SO}$};
\node[below]          at (3, 0)   {$x_{UE}$};

%% y 軸
\draw[thin] (-0.1, 1)   -- (0, 1);
\draw[thin] (-0.1, 2)   -- (0, 2);
\draw[thin] (-0.1, 2.5) -- (0, 2.5);

\node[left, font=\footnotesize] at (0, 1)   {$t_0$};
\node[left]                     at (0, 2)   {$t_{SO}$};
\node[left]                     at (0, 2.5) {$t_{UE}$};

%% =============================================================
%% τ* 注釈（最適課金額：PMC(x_SO)〜SMC(x_SO) の垂直両矢印）
%% =============================================================

\draw[<->, >=stealth, thick]
  (2.2, 2.0) -- (2.2, 3.0)
  node[midway, right, font=\footnotesize] {$\tau^*$};

%% =============================================================
%% 領域ラベル
%% =============================================================

%% 死荷重損失ラベル（三角形の内側）
\node[font=\footnotesize, align=center, violet!70!black]
  at (2.72, 3.15) {死荷重\\[-2pt]損失};

%% =============================================================
%% 均衡点のラベル
%% =============================================================

\node[above left, font=\footnotesize] at (2.2, 3.1)    {SO点};
\node[below right, font=\footnotesize] at (2.8, 2.4) {UE点};

\end{tikzpicture}

%%
%% 模式図のパラメータ（すべて正規化された単位）
%%   需要曲線     D(x)   = 4 - 0.5x         （右下がり）
%%   私的限界費用  PMC(x) = 1 + 0.5x = t_a(x_a)
%%   社会的限界費用 SMC(x) = 1 + x   = d/dx[x * t_a(x_a)]
%%
%%   SO 点 (D = SMC): x_SO = 2,  t_SO = PMC(2) = 2,  価格(含む課金) = 3
%%   UE 点 (D = PMC): x_UE = 3,  t_UE = PMC(3) = 2.5
%%   最適課金 τ* = SMC(x_SO) - PMC(x_SO) = 3 - 2 = 1
%%   自由流旅行時間 t_0 = PMC(0) = SMC(0) = 1
%%   死荷重損失の三角形 頂点: SO点(2,3), UE点(3,2.5), SMC(x_UE)=(3,4)

\begin{tikzpicture}[
  every node/.style={font=\small},
  xscale=1.6,    %% 横方向の拡大率（x 軸 0-4）
  yscale=1.4     %% 縦方向の拡大率（y 軸 0-4.5）
]

%% =============================================================
%% 塗り潰し（背面レイヤー）
%% =============================================================

%% 死荷重損失「紫の三角形」:  SO点 → SMC(x_UE) → UE点 で囲まれる三角形
%% 境界: 上辺 = SMC 曲線 (2,3)→(3,4), 右辺 = 垂直線 (3,4)→(3,2.5),
%%       下辺 = D 曲線 (3,2.5)→(2,3)
\fill[violet!30] (2, 3) -- (3, 4) -- (3, 2.5) -- cycle;

%% =============================================================
%% 曲線（3 本）
%% =============================================================

%% 需要曲線  D(x) = 4 - 0.5x
\draw[very thick, blue!75!black]
  (0, 4) -- (3.6, 2.2)
  node[right] {$D$};

%% 私的限界費用  PMC(x) = t_a(x_a) = 1 + 0.5x
\draw[very thick, green!55!black]
  (0, 1) -- (3.6, 2.8)
  node[right] {$\text{PMC} = t_a(x_a)$};

%% 社会的限界費用  SMC(x) = 1 + x （= d/dx[x * t_a] = t_a + x_a t_a'）
\draw[very thick, red!65!black]
  (0, 1) -- (3.3, 4.3)
  node[right, align=left] {$\text{SMC} = t_a + x_a\, t_a'$};

%% =============================================================
%% 補助破線
%% =============================================================

%% x_SO の垂直破線
\draw[gray, dashed, thin] (2, 0) -- (2, 3);

%% x_UE の垂直破線
\draw[gray, dashed, thin] (3, 0) -- (3, 2.5);

%% t_SO = PMC(x_SO) = 2 の水平破線
\draw[gray, dashed, thin] (0, 2) -- (2, 2);

%% t_UE = PMC(x_UE) = 2.5 の水平破線
\draw[gray, dashed, thin] (0, 2.5) -- (3, 2.5);

%% =============================================================
%% 均衡点（塗り潰し丸）
%% =============================================================

%% SO 点（D と SMC の交点, x=2, y=3）
\fill (2, 3) circle[radius=2.3pt];

%% UE 点（D と PMC の交点, x=3, y=2.5）
\fill (3, 2.5) circle[radius=2.3pt];

%% PMC(x_SO) の点（課金の下端, x=2, y=2）
\fill[gray!60] (2, 2) circle[radius=1.8pt];

%% =============================================================
%% 軸（前面レイヤー）
%% =============================================================

\draw[thick, ->] (-0.15, 0) -- (4.0, 0) node[right] {$x_a$（交通量）};
\draw[thick, ->] (0, -0.15) -- (0, 4.7) node[above] {費用（旅行時間）};

%% 原点ラベル
\node[below left, font=\small] at (0, 0) {O};

%% =============================================================
%% 軸の目盛り・ラベル
%% =============================================================

%% x 軸
\draw[thin] (2, 0) -- (2, -0.1);
\draw[thin] (3, 0) -- (3, -0.1);
\node[below]          at (2, 0)   {$x_{SO}$};
\node[below]          at (3, 0)   {$x_{UE}$};

%% y 軸
\draw[thin] (-0.1, 1)   -- (0, 1);
\draw[thin] (-0.1, 2)   -- (0, 2);
\draw[thin] (-0.1, 2.5) -- (0, 2.5);

\node[left, font=\footnotesize] at (0, 1)   {$t_0$};
\node[left]                     at (0, 2)   {$t_{SO}$};
\node[left]                     at (0, 2.5) {$t_{UE}$};

%% =============================================================
%% τ* 注釈（最適課金額：PMC(x_SO)〜SMC(x_SO) の垂直両矢印）
%% =============================================================

\draw[<->, >=stealth, thick]
  (2.2, 2.0) -- (2.2, 3.0)
  node[midway, right, font=\footnotesize] {$\tau^*$};

%% =============================================================
%% 領域ラベル
%% =============================================================

%% 死荷重損失ラベル（三角形の内側）
\node[font=\footnotesize, align=center, violet!70!black]
  at (2.72, 3.15) {死荷重\\[-2pt]損失};

%% =============================================================
%% 均衡点のラベル
%% =============================================================

\node[above left, font=\footnotesize] at (2.2, 3.1)    {SO点};
\node[below right, font=\footnotesize] at (2.8, 2.4) {UE点};

\end{tikzpicture}

  \caption{余剰分析：UE と SO の比較．
           横軸は交通量 $x_a$，縦軸は費用（旅行時間・支払意思額）．
           $D$ は需要曲線（支払意思額），PMC は私的限界費用（旅行時間 $t_a(x_a)$），
           SMC は社会的限界費用 $\dfrac{d(x_a t_a)}{dx_a}$．
           SO 点（$x_{SO}$）では $D = \text{SMC}$；UE 点（$x_{UE}$）では $D = \text{PMC}$．
           紫の三角形（死荷重損失）は UE 配分によって生じる社会的余剰の損失を表す．
           $\tau^*$ は最適混雑課金額（外部コスト）．}
  \label{fig:surplus-analysis}
\end{figure}

\subsubsection*{社会的余剰の定義と最大化}

需要曲線 $D(x_a)$ は，交通量 $x_a$ 台目の移動がもたらす社会的便益（\textbf{支払意思額}）を表す．
社会的余剰は，総便益から総社会的費用を差し引いた量として定義される：
\begin{equation}
  W(x_a)
  = \int_0^{x_a} D(v)\,dv
    - \int_0^{x_a} \text{SMC}(v)\,dv
  = \int_0^{x_a} \bigl[D(v) - \text{SMC}(v)\bigr]\,dv
  \label{eq:social-welfare}
\end{equation}
この $W$ を最大化する条件は：
\begin{equation}
  \frac{d W}{d x_a} = D(x_a) - \text{SMC}(x_a) = 0
  \quad \Longleftrightarrow \quad
  D(x_a) = \text{SMC}(x_a)
  \label{eq:so-condition}
\end{equation}
すなわち，需要曲線と SMC 曲線が交わる点が社会的余剰を最大化する配分——
これが SO 配分 $x_{SO}$ にほかならない（図\ref{fig:surplus-analysis} の「SO 点」）．

\subsubsection*{UE 配分における死荷重損失}

UE では各利用者が PMC のみを参照するため，
$D(x_a) = \text{PMC}(x_a)$ となる点 $x_{UE}$ で均衡する．
$x_{SO} < x_{UE}$ であるため，交通量は社会的最適水準を超えて過剰に供給される．

この過剰供給によって生じる社会的余剰の損失（\textbf{死荷重損失}，Deadweight Loss）は：
\begin{equation}
  \Delta W
  = \int_{x_{SO}}^{x_{UE}} \bigl[\text{SMC}(x) - D(x)\bigr]\,dx
  \quad > 0
  \label{eq:dwl}
\end{equation}
すなわち，$x_{SO}$ から $x_{UE}$ の範囲で SMC が $D$ を上回る部分の面積（図\ref{fig:surplus-analysis} の
紫の三角形）に等しい．

\subsubsection*{混雑課金による死荷重損失の解消}

最適課金 $\tau_a^* = x_a^{SO}\, t_a'(x_a^{SO})$（式\eqref{eq:optimal-toll}）を
各リンクに設定すると，利用者が直面する一般化費用が PMC から SMC へシフトし，
UE 点が $x_{UE}$ から $x_{SO}$ に移動する．
このとき：

\begin{itemize}
  \item 死荷重損失 $\Delta W$（紫の三角形）が解消される．
  \item 課金収入 $\tau_a^* \times x_a^{SO}$ は政府が徴収する．
        この収入を利用者に一括再配分（lump-sum transfer）することで，
        パレート改善（全員の厚生が改善）が原理的に可能である．
  \item 実際の旅行時間 $t_{SO} = t_a(x_a^{SO})$ は，
        課金前の UE における旅行時間 $t_{UE} = t_a(x_a^{UE})$ より短くなる
        （$x_{SO} < x_{UE}$ かつ $t_a$ 単調増加より）．
\end{itemize}

\begin{tcolorbox}[
  colback=blue!4!white, colframe=blue!50!black,
  title={\small まとめ：UE・SO・混雑課金の関係},
  breakable, fonttitle=\bfseries\small
]
  \begin{itemize}\setlength{\itemsep}{2pt}
    \item UE：各利用者が PMC を最小化 $\Rightarrow$ 外部コストを無視 $\Rightarrow$ 過剰交通量
    \item SO：SMC を最小化する中央集権的配分 $\Rightarrow$ 社会的余剰最大
    \item 限界費用課金 $\tau_a^* = x_a^{SO}\, t_a'(x_a^{SO})$：
          外部コストを利用者に課すことで，UE 解を SO 解に一致させる
    \item 死荷重損失（紫の三角形）は課金によって回収される社会的便益に対応する
  \end{itemize}
\end{tcolorbox}
